\section{The exact binary certificate and its denominator scale}
\label{sec:exact-certificate}

This appendix records a closed certificate that uses the neighboring-length
relations with their exact tangent probabilities. It gives a useful
comparison with the conserved rounding in Section~\ref{sec:rounding}.
The exponential denominator bound below concerns this exact certificate;
it is not a lower bound for rounded or multiway constructions.

Throughout this appendix, \(K,U,T\), the split counts, and the catalyst
refer only to the exact binary certificate. They are distinct from the
rounded quantities in the main construction.

\subsection{A certificate with fixed rational probabilities}

For a fixed target index \(a\ge2\), use rows \(2\le r\le a\) and
columns \(r\le b\le4r-2\). The affine formula in
\eqref{eq:profile-formula} holds in row \(r\) throughout \(b\ge r-1\):
at its left endpoint this follows from
\(s_r/s_{r-1}=(3r-2)/(3r-3)\).
Set
\begin{equation}\label{eq:probabilities}
 q_{r,b}^{\pm}=\frac{2b+r-1\pm2}{2(2b+r-1)}.
\end{equation}
These positive rational probabilities equal the neighboring values of
\(F\) divided by \(2F(r,b)\).
Put \(y_{r,b}=\log(P(r,b)/F(r,b))\). Weighted arithmetic--geometric
mean applied to \eqref{eq:concavity} yields
\begin{equation}\label{eq:log-concavity}
 y_{r,b}-q^-_{r,b}y_{r,b-1}-q^+_{r,b}y_{r,b+1}\ge0.
\end{equation}
Indeed the normalized concavity inequality is
\[
 \frac{P(r,b)}{F(r,b)}\ge
 q^-_{r,b}\frac{P(r,b-1)}{F(r,b-1)}+
 q^+_{r,b}\frac{P(r,b+1)}{F(r,b+1)}
 \ge \exp\bigl(q^-_{r,b}y_{r,b-1}+q^+_{r,b}y_{r,b+1}\bigr).
\]
We will also use the equality
\(H(q)+\sum q^\pm\log F(r,b\pm1)=\log(2F(r,b))\),
where \(H(q)=-\sum q^\pm\log q^\pm\).
Tripling at \(h=r\) gives
\(y_{r,4r-1}-y_{r,r}\ge0\).

Put \(L=4r-1\) and \(D_r=3r-1\). In the final row \(r=a\), set
\begin{equation}\label{eq:final-weights}
 \lambda_{r,b}=\frac{2(2b+r-1)(L-b)}{3(r-1)D_r},\qquad
 \mu_r=\frac1{r-1}.
\end{equation}
In each lower row \(r<a\), set
\begin{equation}\label{eq:lower-weights}
 \begin{split}
 \lambda_{r,r}&=\frac{2D_r}{3(r-1)},\\
 \lambda_{r,b}&=\frac{4(2b+r-1)(L-b)}{3(r-1)(3r+1)}\quad(b>r),\\
 \mu_r&=\frac{2D_r}{(r-1)(3r+1)}.
 \end{split}
\end{equation}

\begin{proposition}[Rational telescoping]\label{prop:dual}
Weight \eqref{eq:log-concavity} by \(\lambda_{r,b}\) and tripling by
\(\mu_r\). Their sum over all specified rows and columns is exactly
\(y_{a,a}-y_{1,2}\ge0\). All weights and all split weights
\(\lambda q^\pm\), \(\mu/3\) are positive rationals.
\end{proposition}
\begin{proof}
Write \(h_b=2b+r-1\) and \(f_b=\lambda_{r,b}/h_b\).
At an interior column the concavity contribution is
\(h_b(f_b-(f_{b-1}+f_{b+1})/2)\).
In the final row \(f_b=2(L-b)/(3(r-1)D_r)\) is affine.
The left external coefficient is
\(-h_{r-1}f_r/2=-1\). At \(r\), the concavity contribution is
\(r/(r-1)\); subtracting \(\mu_r=1/(r-1)\) leaves \(1\).
The terminal coefficient is
\(-h_Lf_{L-1}/2+\mu_r=0\).
Thus this row leaves \(y_{r,r}-y_{r,r-1}\).

In a lower row, \(f_b\) is affine for \(b>r\) and
\(f_r=2/(3(r-1))\). Its affine extension from \(b>r\) differs at
\(r\) by \(2/(3r+1)\). The second difference therefore leaves
coefficient \(1\) at \(r+1\). The left external coefficient remains
\(-1\); the concavity coefficient at \(r\) is
\(2(3r-1)/((r-1)(3r+1))=\mu_r\), which cancels with \(-\mu_r\), and
the terminal coefficient cancels with \(\mu_r\).
This row leaves \(y_{r,r+1}-y_{r,r-1}\).
Symmetry identifies \(y_{r,r+1}=y_{r+1,r}\), so all row boundaries
cancel except \(y_{a,a}-y_{2,1}=y_{a,a}-y_{1,2}\).
\end{proof}

The probabilities in this certificate are fixed by \(F\), not by the
unknown values of \(P\). This fact is what permits an exact finite
tensor construction.

\subsection{Integer types and their entropy identity}
\label{subsec:types}

Let \(K=K_a\) be the least common multiple of the denominators of all
\(\lambda q^+\), \(\lambda q^-\), and \(\mu/3\), in lowest terms.
For each concavity relation write
\[
 w=K\lambda,\qquad u=K\lambda q^+,\qquad d=K\lambda q^-;
 \qquad w=u+d.
\]
For each tripling relation write \(z=K\mu\); then \(3\mid z\).
Let
\begin{equation}\label{eq:type-totals}
 U=\sum w,\quad T=\sum z,\quad
 I_c=\frac{3a^2+a-4}{2},\quad I_t=a-1,\quad I=\frac{3(a+2)(a-1)}2.
\end{equation}
The sums run over concavity and tripling relations, respectively.
There are \(I_c\) and \(I_t\) such relations.

For a formal symbol \([r,b]=[b,r]\), Proposition~\ref{prop:dual} gives
the exact integer identity
\begin{equation}\label{eq:formal-ledger}
 \sum_{r,b}\bigl(w[r,b]-u[r,b+1]-d[r,b-1]\bigr)
 +\sum_r z\bigl([r,4r-1]-[r,r]\bigr)
 =K[a,a]-K[1,2].
\end{equation}
This is an identity of exponent vectors, not an identity of sums of tensors.
Let \(i_{r,b}\) and \(o_{r,b}\) be the nonnegative source and output
counts on its left, after identifying symmetric pairs. Put
\(c_{r,b}=\min(i_{r,b},o_{r,b})\) and
\begin{equation}\label{eq:catalyst}
 \Cat^{\mathrm{ex}}=\bigotimes_{r\le b}Q_{r,b}^{\otimes c_{r,b}},\qquad
 \Cat^{\mathrm{ex}}_N=(\Cat^{\mathrm{ex}})^{\otimes N}.
\end{equation}
Thus \(\sum c_{r,b}=U+T-K\). The common factors are retained as the
catalyst \(\Cat^{\mathrm{ex}}\); they are not formally canceled from a restriction.

Tightness of each entropy relation at \(F\) also gives
\begin{equation}\label{eq:entropy-identity}
 2^{K-U}3^T\prod_{\text{concavity relations}}
 \frac{w^w}{u^u d^d}=F(a,a)^K.
\end{equation}
To check the identity, multiply
\((2F(r,b))^w=(w^w/(u^ud^d))F(r,b+1)^uF(r,b-1)^d\)
and \(F(r,4r-1)^z=3^zF(r,r)^z\) over the relations.
Use \eqref{eq:formal-ledger} and \(F(1,2)=2\).
This derivation proves the identity over the rationals, without a numerical
entropy calculation.

\subsection{The minimal denominator scale}


\begin{theorem}[Exact normalization]\label{thm:normalization}
The least integer scale in Section~\ref{subsec:types} is
\begin{equation}\label{eq:minimal-scale}
 K_a=\lcm\left(
 \frac{3(a-1)(3a-1)}{\gcd(3a-1,4)},\quad
 \left\{\frac{3t(3t+4)}{\gcd(t,2)^2}:1\le t\le a-2\right\}
 \right).
\end{equation}
In particular,
\begin{equation}\label{eq:scale-divisibility}
 3\lcm(1,\ldots,a-1)\mid 2K_a,
 \qquad K_a\mid3\lcm(1,\ldots,3a-1),
\end{equation}
and
\begin{equation}\label{eq:scale-growth}
 K_a\ge\frac32\lcm(1,\ldots,a-1)\ge\frac{3\,2^{a-2}}a.
\end{equation}
\end{theorem}
\begin{proof}
In the final row put \(t=a-1\), \(D_r=3t+2\), and \(b=a+h\).
All split weights have common denominator \(3tD_r\).
Their integer numerators are \(D_r\), from \(\mu/3\), and
\((D_r+2h\pm2)(D_r-h)\) for \(0\le h\le D_r-1\).
Modulo \(D_r\), the latter products are \(-2h(h\pm1)\), all divisible
by four. The case \(h=1\) with the plus sign is congruent to \(-4\).
Thus the joint gcd of the denominator and all numerators is exactly
\(\gcd(D_r,4)\). This gives the first entry in \eqref{eq:minimal-scale}.

In a lower row put \(t=r-1\) and \(V=3t+4\).
Use common denominator \(3tV\). At \(b=r\), the two split numerators
are \(V^2\) and \(3tV\); the tripling split numerator is \(2(3t+2)\).
For \(b=r+h>r\), the split numerators are
\(2(3t+2+2h\pm2)(3t+2-h)\).
An odd prime dividing both \(V^2\) and \(2(3t+2)\) would divide the
difference between \(V\) and \(3t+2\), namely two, which is impossible.
If \(t\) is odd, \(V^2\) is odd, so the joint gcd is one.
If \(t\) is even, all numerators and the denominator are divisible
by four. When \(t\equiv2\pmod4\), \(V^2\) has exactly two factors
of two; when \(t\equiv0\pmod4\), \(2(3t+2)\) has exactly two.
Thus the joint gcd is exactly four. This proves the remaining entries.

Each row denominator is either \(3\lcm(t,3t+c)\) or half that number,
with \(c=2\) for the final row and \(c=4\) for lower rows.
The integers involved lie between one and \(3a-1\); twice each row
denominator is divisible by \(3t\). Taking least common multiples
proves \eqref{eq:scale-divisibility}.

For completeness, every binomial coefficient \(\binom nh\) divides
\(\lcm(1,\ldots,n)\). In its factorial valuation at a prime \(p\),
each power \(p^j\le n\) contributes zero or one, so its valuation
is at most \(\lfloor\log_p n\rfloor\).
The largest binomial coefficient in row \(n\) is at least
\(2^n/(n+1)\). Apply this with \(n=a-1\) and use
\eqref{eq:scale-divisibility} to obtain \eqref{eq:scale-growth}.
\end{proof}

The scale is minimal for these exact rational split weights. Approximate
types, different weights, or different tensor maps need not obey this
bound. Since \(U\ge K_a\), already the sixfold binary source dimension
\(2^{2NU}\) grows at least doubly exponentially in \(a\) under
this bound. The inequality \(U\ge K_a\) also follows directly from
\eqref{eq:formal-ledger}: the net consumption of \([a,a]\) is \(K_a\),
and all its source occurrences belong to concavity relations.
